\section{Kaggle implementation: logistic probabilities and threshold policies}
\label{sec:kaggle-titanic-logistic}

\begin{kaggleproject}{Titanic --- logistic regression as a probability model}
\kagglesource{https://www.kaggle.com/datasets/yasserh/titanic-dataset}{Titanic Dataset}
\textbf{Question.} What happens to precision and recall when the fitted logistic model is unchanged but the classification threshold is moved?
\end{kaggleproject}

\textbf{Before running.} Predict the qualitative pattern you expect from the experiment before you look at the output or plot.

\begin{labmode}
Stage 3B: attempt the probability-to-threshold experiment yourself first. The threshold sweep is performed on a development validation set so it can legitimately inform the operating policy. A final test set, if one exists in a later project, would remain untouched until that policy is locked. The reference script is a verification aid, not the starting point.
\end{labmode}

\textbf{Part A --- fit the logistic model once and obtain validation scores.}

\lstinputlisting[style=mlpython,firstline=1,lastline=40]{all_experiments/lab09-titanic_logistic.py}

\newpage
\textbf{Part B --- change only the decision threshold.}

\lstinputlisting[style=mlpython,firstline=42,lastline=65]{all_experiments/lab09-titanic_logistic.py}


\begin{keyidea}
The probability model and the operational decision rule are separate objects. \texttt{predict\_proba} gives a score; the threshold converts that score into an action.
\end{keyidea}

\begin{labdeliverable}
Submit one notebook with validation predicted probabilities, a table comparing at least three classification thresholds, precision/recall for each threshold, and a short operational recommendation that keeps the fitted probability model separate from the chosen action threshold.
\end{labdeliverable}
